\documentclass[conference]{IEEEtran}

\usepackage{amsmath,amssymb}
\usepackage{booktabs}
\usepackage{graphicx}
\usepackage{url}
\usepackage[hidelinks]{hyperref}

\begin{document}

\title{Correlation-Based Stint Filtering Biases Tyre\\
Degradation Models: Evidence from Three\\
Seasons of Formula~1 Telemetry}

\author{\IEEEauthorblockN{Saptarshi Nag}
\IEEEauthorblockA{\textit{Independent Research} \\
\texttt{saptarshireeju611@gmail.com}}}

\maketitle

\begin{abstract}
Tyre degradation models for Formula~1 conventionally estimate a per-stint
degradation rate as the slope of fuel-corrected lap time against tyre age, and
retain only stints whose linear fit exceeds a correlation threshold. We show
that this near-universal quality criterion is not a quality criterion at all.
Measured across 3\,892 stints from the 2022--2024 seasons, the
absolute correlation $|r|$ tracks slope \emph{magnitude} ($\rho = 0.53$) far
more strongly than it tracks slope \emph{precision} ($\rho = -0.20$), because a
genuinely flat stint has low correlation by construction however well it is
measured. The resulting acceptance rate rises monotonically from 29.5\% for
near-zero slopes to 98.7\% for large ones, systematically deleting
low-degradation circuits from the training set. Replacing the criterion with
one based on the standard error of the fitted slope raises the between-event
share of target variance from 32.3\% to 49.7\%, the cross-season per-circuit
correlation from 0.18 to 0.64, and corrects every remaining directional error
in a downstream pit-strategy simulator. We further report that a
telemetry-derived energy proxy widely treated as an aerodynamic load term
correlates \emph{negatively} with degradation ($\rho = -0.60$), falsifying the
frictional-work hypothesis in that form, and argue from a variance
decomposition that it is better understood as a measure of straight-line
fraction. All results are accompanied by a measured performance ceiling, which
we argue should be established before any improvement threshold is fixed.
\end{abstract}

\begin{IEEEkeywords}
motorsport analytics, tyre degradation, selection bias, time-series
cross-validation, gradient boosting, Formula~1
\end{IEEEkeywords}

\section{Introduction}

\IEEEPARstart{T}{yre} degradation is the dominant controllable variable in
Formula~1 race strategy. A pit stop costs roughly 20--25\,s of track position,
so the decision of when to stop is governed almost entirely by the rate at
which lap time deteriorates with tyre age. Public telemetry, exposed through
the FastF1 library, has made this problem accessible outside the teams
themselves, and a standard modelling recipe has emerged: aggregate laps into
stints, fit a linear model of fuel-corrected lap time against tyre age, and
treat the resulting slope as the degradation rate.

That recipe contains a filtering step that has, to our knowledge, not been
examined. Because a fitted slope may be noisy, stints are retained only when
the regression achieves some minimum absolute correlation, typically
$|r| \geq 0.3$. The stated justification is fit quality. We show the criterion
does not measure fit quality, and that its failure mode is not random but
systematically biased against precisely the circuits where a strategist most
needs an accurate answer.

The contributions of this paper are:

\begin{enumerate}
\item A quantitative demonstration that correlation-based stint filtering
      selects on slope magnitude rather than estimation precision, with an
      acceptance rate varying by a factor of 3.3 across the range of true
      slopes (Section~\ref{sec:bias}).
\item A replacement criterion based on the standard error of the fitted slope,
      and measurement of its effect on variance structure, cross-season
      transferability and downstream strategy correctness
      (Section~\ref{sec:fix}).
\item Evidence that a commonly used energy proxy, $\overline{v^2}$, correlates
      negatively with degradation and is better interpreted as a track
      geometry descriptor than as a contact-patch work term
      (Section~\ref{sec:energy}).
\item An argument, with worked example, that improvement thresholds should be
      set only after the achievable ceiling has been measured
      (Section~\ref{sec:ceiling}).
\end{enumerate}

\section{Related Work}

Rubber friction is viscoelastic, decomposing into adhesion and hysteresis
components; Persson's multiscale theory obtains friction by integrating
viscoelastic losses over the surface roughness spectrum~\cite{persson2001}.
Wear literature descending from Archard establishes that wear rate is
proportional to frictional work dissipated at the contact
patch~\cite{archard1953}, which motivates including an energy term in any
composite wear index.

Direct application of that physics requires a tyre force model. The Pacejka
Magic Formula~\cite{pacejka2012,bakker1986} is the standard, but its
coefficients are load-dependent and fitted to rig data that is proprietary for
Formula~1 constructions. We therefore exclude it, following the same reasoning
as prior public work, and use curvature-derived proxies computed from position
telemetry.

For the learning component, gradient-boosted trees remain state of the art on
medium-sized tabular data, with inductive biases robust to uninformative
features and irregular target functions~\cite{chen2016,grinsztajn2022}. On
validation, standard $k$-fold cross-validation leaks future information when
observations are serially dependent~\cite{bergmeir2018,bergmeir2012}, requiring
chronological splits. Prior Formula~1 work using FastF1 includes tyre energy
forecasting with gradient boosting and recurrent models~\cite{sulsky2025} and a
state-space treatment separating fuel mass from latent tyre
pace~\cite{statespace2025}; both note that the practical industry standard
remains simple linear degradation models.

Derivative estimation from noisy telemetry uses Savitzky--Golay
filtering~\cite{savitzky1964}, whose original coefficient tables contain
typographical errors corrected subsequently~\cite{steinier1972}; we rely on a
library implementation rather than transcribed coefficients.

\section{Data and Methodology}

\subsection{Collection}

We collect all race sessions from the 2022--2024 Formula~1 seasons, a single
ground-effect regulation era, via FastF1. This yields 68 sessions and 49\,993
laps. The API enforces a limit of 500 requests per hour. We measure the true
cost per session before bulk collection: a race loaded with laps and weather
but without telemetry costs 11 requests, a repeat load of the same session
costs 0, and enabling telemetry adds 2.

The zero-cost repeat is architecturally important. The cached session class
places the cache lookup ahead of the rate-limited transport in the method
resolution order, so a cache hit never reaches the limiter. Preserving the
cache therefore makes all subsequent processing free, and we write each session
to its own file through an atomic rename so that interruption costs one session
rather than the corpus.

\subsection{Target construction}

Lap times are fuel-corrected before any slope is fitted. Fuel remaining on
lap~$N$ of a race of $L$ laps, starting from the regulation maximum
$F_0 = 110$\,kg, is
\begin{equation}
F(N) = F_0 - (N-1)\frac{F_0}{L},
\end{equation}
since the car completes lap~$N$ carrying the fuel it held at the start of that
lap. The corrected time is $t^{*} = t - \kappa F(N)$ with
$\kappa = 0.032$\,s/kg. Omitting this correction produces a target dominated by
fuel burn rather than tyre wear.

Representative laps are selected by a mask excluding in-laps and out-laps,
non-green track status, the first lap of each stint, laps slower than 107\% of
the stint median, and laps without a recorded time. This retains 84.8\% of laps.

For each stint $s$, defined by the tuple (season, round, driver number, stint
number), we fit
\begin{equation}
t^{*}_{s}(a) = \beta_{0,s} + \beta_{1,s}\, a + \varepsilon,
\label{eq:slope}
\end{equation}
where $a$ is tyre age in laps. The degradation rate is $\beta_{1,s}$, in
seconds of lap time lost per lap of tyre age.

\subsection{Features and validation}

Features must be knowable \emph{before} a stint is run, since the intended
application chooses stint length. Stint length and clean-lap count are
therefore excluded as reverse-causal. The retained set comprises tyre age at
stint start, compound as both an absolute Pirelli C-number and a within-weekend
relative ordinal, stint mean air and track temperature, three circuit trait
ratings, scheduled race distance, and constructor identity as a one-hot block.

Validation is chronological and grouped on whole events. Splitting rows rather
than events allows stints from a single race to fall on both sides of a fold
boundary, leaking shared track, weather and safety-car conditions. We split on
the ordered set of unique (season, round) keys and map back to rows, asserting
that no event straddles any boundary. Models train on 2022--2023 and are
evaluated on a held-out 2024 season.

\section{Correlation-Based Filtering is Magnitude Selection}
\label{sec:bias}

The conventional criterion retains stint $s$ when $|r_s| \geq 0.3$, where
$r_s$ is the Pearson correlation of the fit in~\eqref{eq:slope}. If this
measures fit quality, $|r_s|$ should correlate with the precision of
$\hat{\beta}_{1,s}$, that is inversely with its standard error
$\mathrm{se}(\hat{\beta}_{1,s})$, and should be independent of the magnitude
$|\hat{\beta}_{1,s}|$.

We fit all 2\,815 dry stints without any filtering and measure both. The result
is the opposite of the stated intent:
\begin{align}
\rho\!\left(|r|,\ |\hat{\beta}_1|\right) &= +0.53 \\
\rho\!\left(|r|,\ \mathrm{se}(\hat{\beta}_1)\right) &= -0.20
\end{align}

The criterion tracks how large the degradation is roughly three times more
strongly than how well it was measured. The mechanism is elementary: for a fit
with residual scatter $\sigma$ over $n$ laps spanning tyre ages with standard
deviation $s_a$,
\begin{equation}
r \approx \frac{\beta_1 s_a}{\sqrt{\beta_1^2 s_a^2 + \sigma^2}},
\end{equation}
so $r \to 0$ as $\beta_1 \to 0$ regardless of how small $\sigma$ is. A stint
with genuinely no degradation cannot achieve high correlation, however
precisely it is measured.

Table~\ref{tab:passrate} shows the consequence. Acceptance rises monotonically
from 29.5\% to 98.7\% across the range of true slopes.

\begin{table}[t]
\caption{Acceptance rate of the $|r| \geq 0.3$ criterion, by true slope
magnitude. A quality criterion would be flat.}
\label{tab:passrate}
\centering
\begin{tabular}{lrr}
\toprule
$|\hat{\beta}_1|$ (s/lap) & Stints & Accepted \\
\midrule
0.00--0.02 & 505 & \textbf{29.5\%} \\
0.02--0.05 & 653 & 86.2\% \\
0.05--0.10 & 935 & 98.7\% \\
0.10--0.20 & 541 & 99.8\% \\
$>0.20$    & 179 & 99.4\% \\
\bottomrule
\end{tabular}
\end{table}

The bias is not uniform across circuits. Monaco, the canonical
low-degradation venue, has a median absolute slope of 0.0175\,s/lap against
0.0614 elsewhere, but a median slope standard error of 0.0062 against 0.0152
elsewhere. Its stints are measured \emph{more} precisely than average and were
being discarded for being flat.

A second, compounding effect arises from the common practice of excluding
negative slopes as physically implausible. Monaco's true median slope is
$-0.008$\,s/lap. A circuit with no measurable degradation produces slopes
scattered either side of zero, and truncating at zero retains only the positive
half of that noise, biasing the estimate upward exactly where negligible
degradation is the operationally important fact.

\section{A Precision-Based Criterion}
\label{sec:fix}

We replace the correlation threshold with a bound on the standard error of the
fitted slope, retaining stint $s$ when
\begin{equation}
\mathrm{se}(\hat{\beta}_{1,s}) \leq \tau, \qquad \tau = 0.02\ \text{s/lap}.
\end{equation}
The value of $\tau$ is chosen relative to the quantity of interest: a typical
degradation rate is 0.065\,s/lap, so $\tau$ admits stints whose rate is pinned
well enough to distinguish a high-degradation circuit from a low one. This
accepts a well-measured flat stint and rejects a noisy one, which is the
property actually required.

The lower plausibility bound is set to $-0.05$\,s/lap rather than zero,
corresponding to approximately one second of improvement across a
twenty-lap stint. This separates track evolution, which genuinely occurs, from
measurement noise around a true rate of zero.

Table~\ref{tab:effect} reports the effect. Every quantity moves in the same
direction, and the change is substantial rather than marginal.

\begin{table}[t]
\caption{Effect of replacing correlation-based with precision-based filtering.}
\label{tab:effect}
\centering
\begin{tabular}{lrr}
\toprule
Quantity & $|r| \geq 0.3$ & $\mathrm{se} \leq 0.02$ \\
\midrule
Between-event variance share & 32.3\% & \textbf{49.7\%} \\
Oracle MAE ceiling & 27.4\% & \textbf{38.0\%} \\
Cross-season circuit $r$ (23 vs 24) & 0.18 & \textbf{0.64} \\
Holdout $R^2$, gradient boosting & $+0.047$ & \textbf{$+0.097$} \\
Holdout $R^2$, linear baseline & $-0.053$ & \textbf{$+0.044$} \\
Retained stints & 2\,101 & 1\,726 \\
\bottomrule
\end{tabular}
\end{table}

The rise in cross-season correlation from 0.18 to 0.64 is the most consequential
result. Under the correlation criterion, a circuit's degradation in one season
barely predicted the next, which would ordinarily be interpreted as genuine
year-to-year instability arising from compound and car changes. It was instead
an artefact of the filter. The signal was present throughout and was being
removed.

We validate the correction against domain ground truth using a pit-strategy
simulator driven by the resulting rates. Under the correlation criterion the
simulator returned a two-stop strategy for Monaco, which is unambiguously
wrong. Under the precision criterion it returns a one-stop for Monaco, Monza,
Spa, Singapore, Jeddah and Silverstone, and a two-stop for Bahrain, Suzuka,
Barcelona and Hungaroring, all consistent with observed practice.

\section{The Energy Proxy Measures Track Geometry}
\label{sec:energy}

Physics-informed designs include an energy term in the target, motivated by the
proportionality of wear to frictional work~\cite{archard1953}. The standard
telemetry proxy is $\overline{v^2}$, the mean squared speed over a stint,
described as an aerodynamic load term.

We compute this from position telemetry over ten races (209 stints). Position
coordinates are reported in units of $10^{-1}$\,m rather than metres; since
curvature scales as the inverse of length, neglecting this yields lateral
accelerations an order of magnitude too small. After correction, the median
peak lateral acceleration is 47.2\,m/s$^2$ (4.8\,g) and median peak combined
acceleration 4.85\,g, both consistent with Formula~1 vehicle dynamics.

The measured correlation with degradation is strong and negative:
$\rho = -0.596$ ($p = 1.9 \times 10^{-21}$). This falsifies the frictional-work
hypothesis in the form stated, which predicts a positive relationship. A
telemetry-free approximation using only speed-trap readings yields
$\rho \approx 0$, so the cheaper substitute produces a false negative and does
not permit the hypothesis to be tested at all.

A variance decomposition explains the sign. 92.6\% of the variance of
$\overline{v^2}$ lies between circuits rather than between stints, so it is a
track descriptor, not a per-stint quantity. Circuits with high mean speed spend
proportionally less of the lap cornering, and cornering rather than speed loads
the tyre. Barcelona, at 196\,km/h mean, averages $+0.117$\,s/lap of degradation;
Jeddah, at 230\,km/h, averages $+0.019$. The circuit-level correlation is
$\rho = -0.717$ ($n=9$).

Compared against hand-assigned circuit severity ratings, $\overline{v^2}$
achieves $\rho = -0.596$ where a subjective composite rating achieves
$+0.578$: comparable strength, opposite sign. We therefore interpret the proxy
as an automatically measured circuit severity descriptor rather than a
contact-patch work term, and suggest that the name \emph{aerodynamic load
proxy} is misleading.

\section{Measure the Ceiling Before Setting the Threshold}
\label{sec:ceiling}

Our pre-registered acceptance criterion required the tuned model to improve on
a linear baseline by at least 15\% in held-out mean absolute error. It does
not: the improvement is $-0.2\%$.

That figure is uninterpretable without a ceiling. We compute an oracle
predictor that assigns each stint the mean degradation of its own event and
compound, which requires knowing the answer in advance and therefore bounds any
model using event-level features. It improves on the global mean by 38.0\%. The
15\% threshold thus demanded 40\% of all achievable headroom, a requirement
fixed before anyone established what was achievable.

The reason for the low ceiling is structural: only 49.7\% of degradation
variance lies between events. The remainder separates stints within a single
race, reflecting traffic, fuel management and driver instruction, none of which
is recoverable from circuit-level features.

\begin{table}[t]
\caption{Held-out performance, 2022--2023 training, 2024 evaluation.}
\label{tab:results}
\centering
\begin{tabular}{lrr}
\toprule
Model & MAE (s/lap) & $R^2$ \\
\midrule
Mean predictor & 0.03539 & --- \\
Linear baseline & \textbf{0.03337} & $+0.044$ \\
Gradient boosting (tuned) & 0.03345 & \textbf{$+0.097$} \\
\bottomrule
\end{tabular}
\end{table}

Table~\ref{tab:results} shows the tuned tree ensemble level with the linear
baseline on absolute error while achieving more than double its $R^2$. Once the
target's construction bias is removed, the relationship between circuit
severity and degradation is close to linear, and the additional flexibility of
boosting contributes only in the tails.

\subsection{Fit adequacy}

We verify that this result reflects the problem rather than a mis-specified
model, via a capacity sweep (Table~\ref{tab:capacity}). Held-out error traces a
clear minimum at 71 trees of depth 5. Below it, training and held-out error are
both high and nearly equal, the signature of insufficient capacity; above it,
training error falls to 0.0145 while held-out error rises, the signature of
memorisation. The selected model sits at the interior optimum. The
training-to-held-out gap is $+14.7\%$, and the cross-validation estimate lies
within 5.0\% of held-out performance, indicating that the grouped chronological
scheme is an honest predictor of generalisation.

\begin{table}[t]
\caption{Capacity sweep. The optimum is interior, bracketed by under- and
over-fitting.}
\label{tab:capacity}
\centering
\begin{tabular}{rrrrr}
\toprule
Trees & Depth & Train MAE & Held-out MAE & Gap \\
\midrule
20   & 2 & 0.03527 & 0.03496 & $-0.9\%$ \\
50   & 3 & 0.03186 & 0.03388 & $+6.4\%$ \\
\textbf{71} & \textbf{5} & \textbf{0.02806} & \textbf{0.03334} & $+18.8\%$ \\
150  & 5 & 0.02450 & 0.03379 & $+37.9\%$ \\
400  & 4 & 0.02223 & 0.03518 & $+58.2\%$ \\
1000 & 8 & 0.01451 & 0.03544 & $+144.3\%$ \\
\bottomrule
\end{tabular}
\end{table}

Split-half reliability of the target, obtained by fitting~\eqref{eq:slope}
separately on odd and even clean laps of each stint, is 0.92 (Pearson). The
target is therefore measured reliably; the difficulty lies in transferring
across events, not in measuring within them.

\section{Discussion}

\subsection{Threats to validity}

The precision threshold $\tau$ is itself a filter on the target and cannot be
entirely free of selection effects. We chose it by reference to the magnitude
of a typical degradation rate rather than by optimising held-out performance,
and we note that a stricter correlation threshold was tested, produced a better
$R^2$, and was rejected on the grounds that its justification would have been
its effect on the metric.

Retention falls from 59.7\% to 44.3\% under the precision criterion. The trade
is favourable on every measured axis, but it is a trade.

Residuals correlate $-0.26$ with position through the held-out season
($p = 2.4\times10^{-11}$), indicating that the model is well fitted but not
uniformly calibrated across a season. We report this rather than correct it,
as any correction would consume the held-out set.

The energy-proxy analysis rests on ten races and nine circuits. The
stint-level correlation is robust ($n=209$, $p < 10^{-20}$), but the
circuit-level figure ($n=9$) is indicative only.

\subsection{Practical implication}

Because per-circuit degradation transfers across seasons at $r = 0.64$--$0.74$
once the bias is removed, while the model captures only 14.4\% of achievable
headroom, a strategy system should prefer directly observed per-circuit rates
where they exist and use a model only to fill gaps. We adopt this ordering, and
record for each estimate whether it is observed or modelled.

\section{Conclusion}

A correlation threshold applied to a regression slope selects for large slopes,
not for well-estimated ones. In tyre degradation modelling this systematically
removes low-degradation circuits, and we have shown it suppresses cross-season
transferability by a factor of three and produces demonstrably wrong strategy
recommendations. Replacing it with a bound on the slope standard error is a
one-line change that recovers the signal.

We further find that the standard $\overline{v^2}$ energy proxy correlates
negatively with degradation and is more accurately described as a measure of
track straight-line fraction, and we argue that improvement thresholds are not
interpretable until the achievable ceiling has been measured.

The broader methodological point is that each of these errors was embedded in
a design that was physically motivated, properly cited, and internally
coherent. None was detectable by inspection; each required measuring the
quantity that the design assumed.

\section*{Reproducibility}

All results are reproducible from the accompanying implementation. Collection
is resumable and cached; the complete pipeline from cold API to trained model
runs in approximately ten minutes of network time and minutes of computation on
a single CPU core. Tuned hyperparameters are persisted and version-checked
against the feature set.

\begin{thebibliography}{00}

\bibitem{persson2001} B.~N.~J. Persson, ``Theory of rubber friction and contact
mechanics,'' \emph{J. Chem. Phys.}, vol.~115, no.~8, pp.~3840--3861, 2001.

\bibitem{archard1953} J.~F. Archard, ``Contact and rubbing of flat surfaces,''
\emph{J. Appl. Phys.}, vol.~24, no.~8, pp.~981--988, 1953.

\bibitem{pacejka2012} H.~B. Pacejka, \emph{Tire and Vehicle Dynamics}, 3rd ed.
Oxford, U.K.: Elsevier, 2012.

\bibitem{bakker1986} E.~Bakker, L.~Nyborg, and H.~B. Pacejka, ``Tyre modelling
for use in vehicle dynamics studies,'' SAE Technical Paper 870421, 1986.

\bibitem{chen2016} T.~Chen and C.~Guestrin, ``XGBoost: A scalable tree boosting
system,'' in \emph{Proc. 22nd ACM SIGKDD}, 2016, pp.~785--794.

\bibitem{grinsztajn2022} L.~Grinsztajn, E.~Oyallon, and G.~Varoquaux, ``Why do
tree-based models still outperform deep learning on typical tabular data?'' in
\emph{Adv. Neural Inf. Process. Syst.}, vol.~35, 2022, pp.~507--520.

\bibitem{bergmeir2018} C.~Bergmeir, R.~J. Hyndman, and B.~Koo, ``A note on the
validity of cross-validation for evaluating autoregressive time series
prediction,'' \emph{Comput. Statist. Data Anal.}, vol.~120, pp.~70--83, 2018.

\bibitem{bergmeir2012} C.~Bergmeir and J.~M. Ben\'{i}tez, ``On the use of
cross-validation for time series predictor evaluation,'' \emph{Inf. Sci.},
vol.~191, pp.~192--213, 2012.

\bibitem{savitzky1964} A.~Savitzky and M.~J.~E. Golay, ``Smoothing and
differentiation of data by simplified least squares procedures,'' \emph{Anal.
Chem.}, vol.~36, no.~8, pp.~1627--1639, 1964.

\bibitem{steinier1972} J.~Steinier, Y.~Termonia, and J.~Deltour, ``Comments on
smoothing and differentiation of data by simplified least square procedure,''
\emph{Anal. Chem.}, vol.~44, no.~11, pp.~1906--1909, 1972.

\bibitem{lundberg2017} S.~M. Lundberg and S.-I. Lee, ``A unified approach to
interpreting model predictions,'' in \emph{Adv. Neural Inf. Process. Syst.},
vol.~30, 2017, pp.~4765--4774.

\bibitem{sulsky2025} ``Explainable time series prediction of tyre energy in
Formula One race strategy,'' arXiv:2501.04067, 2025.

\bibitem{statespace2025} ``A state-space approach to modeling tire degradation
in Formula 1 racing,'' arXiv:2512.00640, 2025.

\end{thebibliography}

\end{document}
